\subsubsection*{1.1.6.}
设 $D$ 是一个圆盘，即与标准圆盘
\[
\{z\in\mathbb C:|z|<1\}.
\]
同胚的拓扑空间。设 $0\in D$，并令 $D^*=D-\{0\}$。假设 $D$ 带有解析流形结构，因而有定向。基本群 $\pi_1(D^*)$ 于是典范同构于 $\mathbb Z$。对以 $0$ 为中心的标准圆盘及 $t\in D^*$，$\pi_1(D^*,t)$ 的正生成元由闭路
\[
x\longmapsto \exp(2\pi i x)t.
\]
表示。

\subsubsection*{1.1.7.}
设 $z:D\longrightarrow\mathbb C$ 是一个一致化参数，即 $(dz)_0\ne0$，并设 $\widetilde D^*$ 是 $D^*$ 的一个万有覆盖。用 $\log z$ 表示 $\widetilde D^*$ 上任意满足
\[
\exp(\log z)=z.
\]
的函数。对偶 $(\widetilde D^*,\log z)$ 在唯一同构意义下唯一。

对 $\alpha\in\mathbb C$，令
\[
z^\alpha=\exp(\alpha\log z).
\]
设 $T$ 是 $\pi_1(D^*)$ 的正生成元，按 1.1.5 作用于 $\widetilde D^*$，并通过搬运结构作用于 $\widetilde D^*$ 上的函数。则
\begin{equation}\tag{1.1.7.1}
T\log z=\log z+2\pi i,
\end{equation}
且
\begin{equation}\tag{1.1.7.2}
Tz^\alpha=\exp(2\pi i\alpha)z^\alpha.
\end{equation}

\subsubsection*{1.1.8.}
设 $D$ 是圆盘，$0\in D$，并令 $D^*=D-\{0\}$。若 $D'\subset D$ 是包含 $0$ 的较小圆盘，则限制到 $D'^{*}$ 给出从 $D^*$ 上局部常值层的范畴到 $D'^{*}$ 上局部常值层的范畴的等价。事实上，对 $t\in D'^{*}$，
\[
\pi_1(D'^{*},t)\xrightarrow{\ \sim\ }\pi_1(D^*,t).
\]

现在假设 $D$ 带有可微流形结构。设 $T$ 是 $D$ 在 $0$ 处的切空间，并令 $T^*=T-\{0\}$。对上述 $D'\subset D$ 以及满足 $f(0)=0$ 和 $(df)_0=\mathrm{id}$ 的嵌入 $f:D'\hookrightarrow T$，复合
\[
(\text{限制到 }D')^{-1}\circ f^*
\]
是从 $T^*$ 上局部常值层的范畴到 $D^*$ 上局部常值层的范畴的等价 $[f]$。在典范同构意义下，它只依赖于 $f$ 在 $0$ 处满足 $(df)_0=\mathrm{id}$ 的芽。这样的两个芽之间的一条道路定义等价之间的一个同构，而同伦的道路定义相同的同构。由于这些芽构成的空间可缩——由 $(1-\lambda)f+\lambda f_0$ 给出向某个固定芽 $f_0$ 的收缩——我们得到从 $T^*$ 上局部常值层到 $D^*$ 上局部常值层的一个在唯一同构意义下良定义的等价。

\subsubsection*{1.1.9.}
设 $D$ 是圆盘，$0\in D$，令 $D^*=D-\{0\}$，并设 $p:\widetilde D^*\longrightarrow D^*$ 是一个万有覆盖。记
\[
p:\widetilde D\longrightarrow D
\]
为在 $\widetilde D^*$ 上添加一个点 $0$ 而得到的 $D$ 上拓扑空间，使得：
\begin{enumerate}[label=(\alph*),leftmargin=2.6em]
\item $\widetilde D^*$ 在 $\widetilde D$ 中开；
\item $p:\widetilde D\longrightarrow D$ 延拓 $p:\widetilde D^*\longrightarrow D^*$，且 $p(0)=0$；
\item 当 $U$ 遍历 $D$ 中 $0$ 的各邻域时，集合 $p^{-1}(U)$ 构成 $\widetilde D$ 中 $0$ 的一个邻域基。
\end{enumerate}

\subsubsection*{1.1.10.}
设 $X$ 是非奇异复代数曲线，$0\in X$，并设 $D$ 是 $X$ 中 $0$ 的一个圆盘邻域。
$D$ 上在 $X$ 的有理函数域 $K(X)$ 上代数的全纯函数在 $0$ 处的芽所构成的环，是 $X$ 在 $0$ 处的局部环 $\mathcal O_0$ 的 Hensel 化 $\mathcal O_0^h$。
设 $\widetilde D^*$ 是 $D^*=D-\{0\}$ 的一个万有覆盖。
$\widetilde D^*$ 上在 $K(X)$ 上代数的亚纯函数在 $0\in\widetilde D$ 处的芽所构成的域，是 $K(X)$ 的一个代数闭包 $\overline K$，因而也是 $\mathcal O_0^h$ 的分式域 $K_0^h$ 的代数闭包。

$\widetilde D^*$ 的 $D$-自同构群 $\pi_1(D^*)$（1.1.5）通过搬运结构作用于 $\overline K$，给出
\begin{equation}\tag{1.1.10.1}
\begin{tikzcd}[column sep=huge, row sep=large]
\pi_1(D^*) \arrow[r] \arrow[d, equal]
  & \operatorname{Gal}(\overline K/K_0^h) \arrow[d, equal]
  & {} \\
\mathbb Z
  & \widehat{\mathbb Z}(1)
  & .
\end{tikzcd}
\end{equation}
由 (1.1.7.2) 可知，此图与在 $\mathbb C$ 上 $\mathbb Z$ 到 $\widehat{\mathbb Z}(1)$ 的自然嵌入相容。此嵌入是 $\mathbb Z\hookrightarrow\widehat{\mathbb Z}$ 与同构
\[
\widehat{\mathbb Z}\xrightarrow{\ \sim\ }
\widehat{\mathbb Z}(1),
\]
的复合，而后一同构是同构
\[
\mathbb Z/n\longrightarrow\mu_n,
\qquad
k\longmapsto\exp\left(\frac{2\pi i}{n}k\right).
\]
的逆极限。

\subsubsection*{1.1.11.}
$\Spec(K_0^h)$ 上的每个层 $F$ 都定义 $D^*$ 上的一个局部常值层。我们使用 1.1.4 的等价，对 $D^*$ 的每个万有覆盖 $\widetilde D^*$，赋予由 1.1.10 中构造的域 $\overline K$ 所确定的 $F_{\overline K}$。
